\documentclass[11pt]{article}

% ===================== PACKAGES =====================
\usepackage[a4paper,margin=1in]{geometry}
\usepackage{amsmath,amssymb}
\usepackage{enumitem}
\usepackage{fancyhdr}
\usepackage{xcolor}

% ===================== HEADER & FOOTER =====================
\pagestyle{fancy}
\fancyhf{}
\lhead{CSE 400: Fundamentals of Probability in Computing}
\rhead{Lecture 17 Scribe}
\cfoot{\thepage}

% ===================== TITLE =====================
\title{
\normalsize School of Engineering and Applied Science (SEAS), Ahmedabad University \\
\vspace{0.2cm}
\textbf{CSE 400: Fundamentals of Probability in Computing}\\
\Large Lecture 17: Joint Random Variables and Joint Distributions
}

\author{}
\date{}

\begin{document}

\maketitle

\vspace{-1.5cm}

\begin{center}
\begin{tabular}{ll}
\textbf{Name:} & {Khushi Paghadar \hspace{2.5in}} \\ [1.5ex]
\textbf{Enrollment No:} & {AU2440131 \hspace{2.5in}} \\ [1.5ex]
\textbf{Course:} & {CSE 400 \hspace{2.8in}} \\ [1.5ex]
\textbf{Topic:} & {Joint Random Variables} \\ [1.5ex]
\end{tabular}
\end{center}

\hrule
\vspace{0.5cm}

\section*{Objective}

This scribe serves as an academic reference for exam preparation for Lecture 17.  
It summarizes the key concepts of joint random variables, joint distributions, and their properties.

\section*{1. Why Study Joint Random Variables?}

In many real-world situations, more than one random variable must be analyzed at the same time.  
Examples include:

\begin{itemize}
\item Smart transportation systems
\item Wireless network performance
\item Weather prediction
\item Drone delivery systems
\item Rolling two dice
\end{itemize}

In such systems, we often want to understand how two variables behave together.

For example:

\[
X = \text{Vehicle Speed}, \quad Y = \text{Traffic Density}
\]

A natural question is whether high vehicle speed can occur when traffic density is also high.  
To answer such questions, we study the \textbf{joint distribution} of the variables.

\section*{2. Joint Random Variables}

Let $X$ and $Y$ be two random variables defined on the same probability space.

Instead of studying them individually, we analyze their \textbf{joint behavior}.  
The pair $(X,Y)$ represents the simultaneous outcome of both variables.

\section*{3. Discrete Case – Joint Probability Mass Function}

For discrete random variables, the joint distribution is defined using the \textbf{Joint Probability Mass Function (PMF)}.

\[
p_{X,Y}(x,y) = P(X = x, Y = y)
\]

This function gives the probability that:

\[
X = x \quad \text{and} \quad Y = y
\]

occur simultaneously.

\subsection*{Properties of Joint PMF}

\textbf{1. Non-negativity}

\[
p_{X,Y}(x,y) \geq 0
\]

\textbf{2. Total Probability}

\[
\sum_x \sum_y p_{X,Y}(x,y) = 1
\]

\section*{4. Joint PMF Table Representation}

A joint PMF can be represented using a two-dimensional probability table.

\[
\begin{array}{c|ccc}
Y \backslash X & x_1 & x_2 & x_3 \\
\hline
y_1 & p(x_1,y_1) & p(x_2,y_1) & p(x_3,y_1) \\
y_2 & p(x_1,y_2) & p(x_2,y_2) & p(x_3,y_2) \\
y_3 & p(x_1,y_3) & p(x_2,y_3) & p(x_3,y_3)
\end{array}
\]

Each cell represents the probability of a specific pair $(x,y)$.

\section*{5. Example: Rolling Two Dice}

Let:

\[
X = \text{value on first die}
\]

\[
Y = \text{value on second die}
\]

Each die can take values:

\[
\{1,2,3,4,5,6\}
\]

Since the dice are fair:

\[
P(X=x,Y=y) = \frac{1}{36}
\]

for all possible pairs.

Example:

\[
P(X+Y=7)
\]

Possible outcomes:

\[
(1,6),(2,5),(3,4),(4,3),(5,2),(6,1)
\]

Thus:

\[
P(X+Y=7)=\frac{6}{36}=\frac{1}{6}
\]

\section*{6. Continuous Case – Joint Probability Density Function}

For continuous random variables, the joint distribution is defined using the \textbf{Joint Probability Density Function (PDF)}.

\[
f_{X,Y}(x,y)
\]

Probability over a region $A$ is given by:

\[
P((X,Y)\in A)=\iint_A f_{X,Y}(x,y)\,dx\,dy
\]

\subsection*{Properties}

\textbf{Non-negativity}

\[
f_{X,Y}(x,y)\ge0
\]

\textbf{Total Probability}

\[
\int_{-\infty}^{\infty}\int_{-\infty}^{\infty} f_{X,Y}(x,y)\,dx\,dy=1
\]

\section*{7. Geometric Interpretation}

In the continuous case:

\begin{itemize}
\item The joint PDF forms a surface above the $xy$-plane.
\item The probability of a region corresponds to the \textbf{volume under the surface}.
\end{itemize}

If $R$ is a region:

\[
P((X,Y)\in R)=\iint_R f_{X,Y}(x,y)\,dx\,dy
\]

\section*{8. Joint Cumulative Distribution Function}

The \textbf{Joint CDF} is defined as:

\[
F_{X,Y}(x,y)=P(X\le x, Y\le y)
\]

It represents the probability that both variables are less than or equal to specific values.

\section*{9. Joint CDF – Continuous Case}

For continuous random variables:

\[
F_{X,Y}(x,y)=\int_{-\infty}^{x}\int_{-\infty}^{y} f_{X,Y}(u,v)\,dv\,du
\]

\section*{10. Joint CDF – Discrete Case}

For discrete random variables:

\[
F_{X,Y}(x,y)=\sum_{u\le x}\sum_{v\le y} p_{X,Y}(u,v)
\]

Thus:

\begin{itemize}
\item Continuous case → double integration
\item Discrete case → double summation
\end{itemize}

\section*{11. Properties of Joint CDF}

\textbf{Bounds}

\[
0 \le F_{X,Y}(x,y) \le 1
\]

\textbf{Monotonicity}

The function is non-decreasing in both arguments.

\textbf{Limits}

\[
\lim_{x\to\infty,\,y\to\infty} F_{X,Y}(x,y)=1
\]

\textbf{Rectangle Probability}

For region

\[
a<X\le b,\quad c<Y\le d
\]

\[
P(a<X\le b, c<Y\le d)=F(b,d)-F(a,d)-F(b,c)+F(a,c)
\]

\section*{12. Marginal Distributions}

Marginal distributions can be obtained from the joint distribution.

\textbf{Discrete case}

\[
P_X(x)=\sum_y p_{X,Y}(x,y)
\]

\[
P_Y(y)=\sum_x p_{X,Y}(x,y)
\]

\textbf{Continuous case}

\[
f_X(x)=\int_{-\infty}^{\infty} f_{X,Y}(x,y)\,dy
\]

\[
f_Y(y)=\int_{-\infty}^{\infty} f_{X,Y}(x,y)\,dx
\]

\section*{13. Key Takeaways}

\begin{enumerate}
\item Joint random variables describe the simultaneous behavior of two variables.
\item Discrete distributions use \textbf{Joint PMF}.
\item Continuous distributions use \textbf{Joint PDF}.
\item Joint probabilities in continuous cases require \textbf{double integrals}.
\item Joint CDF represents cumulative probability of both variables.
\item Marginal distributions are obtained from the joint distribution.
\end{enumerate}

\bigskip
\hrule
\vspace{0.2cm}

\begin{center}
\small \textit{End of Lecture Scribe}
\end{center}

\end{document}